% \begin{savequote}
% \includegraphics[width=6cm]{./images/gabriel-sebastian-aass/cape_town_paramedic.jpg}
% 
% {\sffamily\tiny Bild: Gabriel}
% \end{savequote}
\def\ChapterCode{ALS}
\chapter
[\CO{[ALS]} Advanced Life Support (ALS)]
{Advanced Life Support (ALS)}
\label{chp:ALS}

\ChapterIntro
{%
}
{%
\begin{description}
  \item[Maintainer:] Sebastian Gabriel
  \item[Autoren:] Diverse
  \item[Reviewer:] Standard-Reviewprozess
  \item[Version:] {\VarDocVersionTypeName} ({\VarDocVersionTypeDescription})
  \item[SHA1:] {\DocGitCommit}
\end{description}

{\TxTeilkapitelAASS}
}


\clearpage % TEMP temp clearpage
\Text
{{\TxBeschreibung}}%1 Überschrift
{%
\I{Advanced Life Support} (\I{ALS}) bezeichnet eine Fülle
von erweiterten Maßnahmen zur Therapie des
Herz-Kreislaufstillstands, bzw. im weiteren Sinne von
ähnlichen lebensbedrohlichen Krankheitsbildern. Sie umfassen
unter anderem das Setzen venöser Zugänge, die medikamentöse
Therapie, bestmögliche Atemwegssicherung und Defibrillation.
\Abk{ALS}-Maßnahmen können 
u.\,a. von Sanitätern (ab der Qualifikation \Abk{NKV}) oder
(Not-)Ärzten durchgeführt werden. Der jeweilige Umfang der Maßnahmen
richtet sich nach dem Ausbildungsstand, der eigenen
Erfahrung und dem vorhandenen Material. Sind die
Voraussetzungen für eine Reanimation mittels \Abk{ALS}
nicht gegeben, so muss eine Basisreanimation (wenn
vorhanden mittels \Abk{AED}) durchgeführt werden und weitere
Hilfe angefordert werden. 

Die ERC-Guidelines sind in nach Themen in elf Sektionen unterteilt:
\begin{compactenum}
  \item Kurzdarstellung 
  \cite{Erc-2015-Sektion-01-De,Erc-2015-Section-01}

  \item Basismaßnahmen zur Wiederbelebung Erwachsener 
  und Verwendung automatisierter externer Defibrillatoren 
  \cite{Erc-2015-Sektion-02-De,Erc-2015-Section-02}

  \item Erweiterte Reanimationsmaßnahmen für Erwachsene 
  (\enquote{adult advanced life support}) 
  \cite{Erc-2015-Sektion-03-De,Erc-2015-Section-03}

  \item Kreislaufstillstand in besonderen Situationen 
  \cite{Erc-2015-Sektion-04-De,Erc-2015-Section-04}

  \item Postreanimationsbehandlung 
  \cite{Erc-2015-Sektion-05-De,Erc-2015-Section-05}

  \item Lebensrettende Maßnahmen bei Kindern 
  (\enquote{paediatric life support})
  \cite{Erc-2015-Sektion-06-De,Erc-2015-Section-06}

  \item Die Versorgung 
  und Reanimation des Neugeborenen 
  \cite{Erc-2015-Sektion-07-De,Erc-2015-Section-07}

  \item Das initiale Management des akuten Koronarsyndroms
  \cite{Erc-2015-Sektion-08-De,Erc-2015-Section-08}

  \item Erste Hilfe 
  \cite{Erc-2015-Sektion-09-De,Erc-2015-Section-09}

  \item Ausbildung 
  und Implementierung der Reanimation 
  \cite{Erc-2015-Sektion-10-De,Erc-2015-Section-10}

  \item Ethik der Reanimation 
  und Entscheidungen am Lebensende 
  \cite{Erc-2015-Sektion-11-De,Erc-2015-Section-11}
\end{compactenum}
}%2 Text


\section{\HeadingKeyword{Leitlinien}}

\TextSlide{ERC}
{%
\label{topic:erc}
\index{Fachgesellschaft!ERC}
Eine wichtige Fachgesellschaft, welche sich dem Thema der
Reanimation und damit verwandter Gebiete widmet, ist das
\T[ERC|see{European Resuscitation Council}]{ERC}
(\I{European Resuscitation Council}, \cite{www:Erc}). 
%
Bekannt ist das
\Abk{ERC} vor allem durch seine Leitlinien
(\I{Guidelines}\footnote{\Lang{engl.} \emph{Guideline}:
Leit-, Richtlinie, Empfehlung}) zur Wiederbelebung, welche
alle fünf Jahre aktualisiert werden 
\cite{%
Erc-2015-Section-01,
Erc-2015-Section-02,
Erc-2015-Section-03,
Erc-2015-Section-04,
Erc-2015-Section-05,
Erc-2015-Section-06,
Erc-2015-Section-07,
Erc-2015-Section-08,
Erc-2015-Section-09,
Erc-2015-Section-10,
Erc-2015-Section-11,
}. 
%
Die Guidelines
definieren \E{nicht} den einzig richtigen Weg wie eine
Reanimation ablaufen muss. Vielmehr sind sie als weithin
akzeptierte Ansicht zu
verstehen, wie sie sicher und effektiv durchgeführt werden
kann. \cite[p.1219]{Erc2010Sec1}

Jeder neuen Auflage geht ein aufwendiger Forschungs- und
Diskussionsprozess voraus: Es werden hunderte Studien mit
vielen verschiedenen Fragestellungen an vielen
Forschungszentren auf der ganzen Welt durchgeführt, die
Ergebnisse ausgewertet und anschließend diskutiert.
Die derzeit aktuelle Version
der ERC-Guidelines kam im Jahr \EE{2015} heraus
(\I{ERC 2015}), die nächste Aktualisierung ist im Jahr
2020 zu erwarten. 

Die ERC-Leitlinien werden im Journal 
\I[Resuscitation!Journal]{Resuscitation} 
publiziert 
und sind sowohl auf der Homepage des Journals 
als auch der ERC kostenlos herunterladbar \cite[www:ResuscitationJournal,www:Erc].
%
Eine deutsche Übersetzung 
ist auf den Homepages 
der deutschen 
und der österreichischen Fachgesellschaften 
verfügbar\footnote{%
\I{GRC}: \I{German Resuscitation Council}; 
\I{ARC}: \I{Austrian Resuscitation Council}%
} %footnote
\cite{Erc-2015-Sektion-01-De,
Erc-2015-Sektion-02-De,
Erc-2015-Sektion-03-De,
Erc-2015-Sektion-04-De,
Erc-2015-Sektion-05-De,
Erc-2015-Sektion-06-De,
Erc-2015-Sektion-07-De,
Erc-2015-Sektion-08-De,
Erc-2015-Sektion-09-De,
Erc-2015-Sektion-10-De,
Erc-2015-Sektion-11-De,}%
.

}{%
\begin{itemize}
   \item Leitlinien zur Reanimation
   \item Aktualisierung alle 5 Jahre
   \item Aktuelle Version: 2015
\end{itemize}

}{}{}{}





\Text{Neu oder alt?}
{%
Da alle fünf Jahre eine neue Version der Richtlinien
erscheint, kommt es gerade in der
Übergangszeit\footnote{Als Übergangszeit versteht man den
Zeitraum zwischen Veröffentlichung einer neuen Version der
Leitlinie und Implementierung durch die jeweilige
Organisation.} regelmäßig zu Verwirrungen, welche Versionen
oder Varianten in der Praxis angewendet werden sollen.

\EE{Grundsätzlich muss das gesamte Personal einer Organisation
\E{das von der eigenen Organisation empfohlene} oder
angepasste Schema kompetent im Team anwenden können.} Das
behandelnde Fachpersonal kann jedoch davon im Rahmen der
\E{Eigenverantwortung} abweichen, sofern die dazu
notwendigen Voraussetzungen gegeben sind. Dabei ist ganz
besonders zu beachten, dass das \E{gesamte} Team mit der
anderen oder neueren Version vertraut und die entsprechende
Routine besitzen muss. Darüber hinaus müssen auch
z.\,B. technische Gegebenheiten bedacht werden
(Programmierung der AEDs, etc.). I.\,d.\,R. trifft der
Teamleiter als Letztverantworlicher die
Entscheidung, er muss daher auch genau über die jeweilige
Qualifikation und Routine seiner Teammitglieder Bescheid
wissen. Ist die Durchführung einer anderen Variante als die
von der Organisation empfohlene aufgrund der Ausbildung, des
Trainings und der Routine von Team-Mitgliedern oder der
vorhandenen Ausstattung fragwürdig, so ist im Zweifel im Sinne eines
reibungslosen Ablaufes des empfohlene Schema zu verwenden.
\E{Die Publikation einer neuen Ausgabe der Guidelines
bedeutet nicht, das die bisherigen Methoden unsicher oder
ineffektiv sind \cite[p.1219]{Erc2010Sec1}.}

% \enquote{The guidelines [..] do not define the only way that
% resuscitation can be delivered; they merely represent a
% widely accepted view of how resuscitation should be
% undertaken both safely and effectively. The publication of
% new and revised treatment recommendations does not imply
% that current clinical care is either unsafe or ineffective.}

}
%{%
%\begin{itemize}
%   \item
%\end{itemize}
%
%}{}{}{}


% \EE{Sonderfall: VF/VT unter bestimmten Bedingungen}:
% Dies betrifft
% Patienten mit Kammerflimmern oder ventrikulärer
% Tachykardie \Z{während einer Herzkatheteruntersuchung oder
% früh-postoperativ nach einem herzchirugischen Eingriff bei
% einem beobachteten Kreislaufstillstand, wenn der Patient
% bereits an einem manuellen Defibrillator angeschlossen
% ist: Hier dürfen 3 Schocks hintereinander gegeben
% werden} Unter Bedachtnahme des Abschnitts bezüglich der
% Elektrotherapie gilt dies auch für Patienten, welche einen
% beobachteten Kreislaufstillstand erleiden und zu dem
% Zeitpunkt bereits an einem Defibrillator angeschlossen
% sind.

\clearpage
\section{\HeadingKeyword{Säulen} der Reanimation}

\subsection{\HeadingKeyword{Basisreanimation}}


\Text
{Einleitung}%1 Überschrift
{%
Das unter \CO{[EHI]} (\prettyref{chp:EHI}) und \CO{[BLS]}
(\prettyref{chp:BLS}) Gesagte gilt auch beim Advanced
Life Support. 
%
Auch hier muss besonders darauf geachtet
werden, 
die Herzdruckmassage so wenig wie möglich zu pausieren 
und mit möglichst optimaler Qualität
durchzuführen. 
%
Unterbrechungen sind nur statthaft, um
spezielle Interventionen zu ermöglichen.
}%2 Text

%\Text{Neuerungen 2010}{%
%Die Empfehlung des
%\enquote*{Pflicht-Zyklus vor AED} beim unbeobachteten
%Kreislaufstillstand ist gestrichen
%worden, kann aber weiterhin angewendet werden.
%Während der Ladephase des
%Defibrillators soll die Herzdruckmassage (\I{HDM})
%fortgeführt werden um die HDM-Pausen zu reduzieren,
%d.\,h. nach der Analyse wird weiter massiert, bis der Defi bereit zur
%Schockabgabe ist.  Bei
%vielen AED-Geräten kann es jedoch zu
%\EE{Problemen} kommen, sie interpretieren die
%reanimationsbedingten Störungen der EKG-Ableitung als
%Kontaktverlust bzw. Bewegungsartefakt und brechen den
%Ladevorgang ab.\footnote{Davon betroffen sind derzeit z.\,B. die
%Geräte der Schiller Fred{\texttrademark}-Serie, der
%Defigard{\texttrademark} 1002 und der
%LifePak{\texttrademark} 500. Bezüglich anderer Geräte liegen
%der Redaktion derzeit noch keine Informationen vor. Es muss
%jedoch davon ausgegangen werden, dass auch andere Modelle
%ein derartiges Verhalten zeigen können. \E{(Angaben ohne
%Gewähr. Stand: 2011-01-19)}} Dadurch wird die Schockabgabe
%verhindert. Daher kann diese Methode nicht generell
%empfohlen werden, sondern darf nur durchgeführt werden, wenn
%der Hersteller des Gerätes dies freigegeben hat.
%\Commentary[HDM während Defi lädt / Doppelte Warnung]{Eine
%Warnung betreffend der HDM während der Ladephase des Defis
%findet sich bereits im \CO{[BLS]}-Teil
%(\prettyref{chp:BLS}). Aufgrund der Tragweite wurde sie im
%\CO{[ALS]}-Teil praktisch
%wortwörtlich wiederholt.}
%}{}{}{}{}
%
%\begin{Danger}
%  \item Manche Defibrillatoren brechen den Ladevorgang ab,
%  wenn währenddessen eine Herzdruckmassage durchgeführt wird.
%\end{Danger}
%\begin{Synopsis}
%  \SynopsisItem{Während des Ladevorganges des Defibrillators darf nur
%  bei Verwendung von \E{manuellen} Geräten eine
%  Herzdruckmassage durchgeführt werden.}
%\end{Synopsis}

\subsection{\HeadingKeyword{Elektrotherapie}}



\TextSlide{\TxBeschreibung}
{%
Die Elektrotherapie hat das Ziel, 
bei Vorliegen einer Herzrhythmusstörung einen normalen Sinusrhythmus mittels elektrischer Impulse wiederherzustellen. 
%
Man nennt diesen Vorgang  \T[Kardioversion!elektrische]{elektrische Kardioversion}%
\ZFootnote{Eine Kardioversion, 
also die Überführung in einen anderen Rhythmus, 
kann auch mittels Medikamenten erfolgen.}% ZFootnote
. 
%
Ein Sonderfall ist die \T{Defibrillation}, sie ist eine Kardioversion bei Kammerflimmern \Z{(Flimmern = Fibrillation)}. 
%
Im Rahmen der Reanimation werden die Begriffe Kardioversion und Defibrillation jedoch umgangssprachlich oft synonym verwendet.%
\ZFootnote{Die elektrische Kardioversion wird auch bei nicht akut lebensbedrohlichen Rhythmusstörungen, 
wie zum Beispiel beim Vorhofflimmern, 
eingesetzt.
%
Die Anwendung erfolgt dann in Narkose bzw. Sedierung 
und es werden geringere Energien eingesetzt. %
}% ZFootnote

Grundsätzlich gilt bezüglich der Defibrillation das unter  \CO{[BLS]} (\FullPrettyRef{topic:defibrillation-bls}) Geschriebene. 
%
Oft ist eine Störung des Reizleitungssystems des Herzens bzw. des elektrischen Herzrhythmus die
Ursache des Kreislaufstillstands. 
Die Elektrotherapie im Sinne der \I{Defibrillation} soll die normale elektrische Herztätigkeit wiederherstellen. 
}
{%
\begin{itemize}
  \item Kardioversion: Überführen in einen anderen Rhythmus
  \item Defibrillation
  \item Siehe \CO{[BLS]} (\FullPrettyRef{topic:defibrillation-bls})
\end{itemize}
}



\TextSlide{Defibrillatoren}{
Geräte für die Defibrillation werden \II{Defibrillator}
genannt. Es gibt folgende Gerätekategorien:
\begin{itemize}
  \item Automatische und halbautomatische Defibrillatoren
  (\T{AED} \I{Automatischer
    externer Defibrillator}, 
  \T{SAED} \I{Semi-automatischer externer Defibrillator}, \enquote*{Halbautomat}):
  %
  Das Gerät analysiert das EKG und
  beurteilt, ob ein Schock abgegeben werden soll. 
  %
  Bei automatischen Defibrillatoren wird die Analysefunktion vom Gerät automatisch gestartet, 
  diese Kategorie kann von Laienhelfern bedient werden. 
  %
  Bei halbautomatischen Geräten wird die Analyse manuell gestartet,
  sie sind für nicht-ärztliches Fachpersonal konzipiert.
  %
  Die Schockauslösung erfolgt bei beiden Arten durch den Anwender. 
  \item Manuelle Defibrillatoren: Das Gerät zeigt das EKG an, der Anwender muss
  entscheiden, ob ein Schock abgegeben werden soll. 
  Manuelle
  Defibrillatoren dürfen nur durch entsprechend ausgebildetes und befugtes 
  (i.\,d.\,R. ärztliches)  
  Personal angewendet werden. 
  Viele neuere Geräte verfügen auch über einen
  \enquote{Halbautomaten-Modus} ((S)AED-Modus), in welchem das Gerät als
  (S)AED verwendet werden kann.
  \item Automatische implantierbare Defibrillatoren
  können Patienten operativ -- ähnlich einem Herzschrittmacher --
  implantiert werden (\Lang{Abkz.} \I{ICD}: \I{Implantable Cardiac Device}).
\end{itemize}

Grundsätzlich kann eine Reanimation auch beim ALS mittels AED durchgeführt werden.
%
Es empfiehlt sich jedoch, 
sofern möglich, 
in den manuellen Modus zu wechseln, 
wenn geeignetes Personal zur Verfügung steht.
%
Dadurch kann die Analysezeit oft reduziert werden 
und auf bestimmte Situationen besser reagiert werden 
(Intubation, \ldots).
}{
\begin{itemize}
  \item Halbautomatische Defibrillatoren (AED)
  \item Manuelle Defibrillatoren
  \item \Z{Automatische implantierbare Defibrillatoren}
\end{itemize}
}{}{}{}



\TextSlide{Elektrische Herzrhythmen}{
Die Defibrillation ist nicht bei jeder Störung des
Herzrhythmus sinnvoll. Man unterscheidet daher schockbare
und nicht-schockbare Rhythmen:
\begin{itemize}
  \item \EE{Schockbare Rhythmen}
\begin{itemize}
    \item \II{Kammerflimmern}: Unkoordnierte, \enquote{wirre}
    elektrische Aktivität
    \item \II{Pulslose Ventrikuläre Tachykardie} (\I{PVT}):
    Zu schneller Herzrhythmus, das Herz kann sich nicht
    ausreichend füllen, daher fehlt die Pumpleistung
  \end{itemize}
  \item \EE{Nicht-schockbare Rhythmen}
  \begin{itemize}
    \item \II{Asystolie}: Keine elektrische Aktivität
    \item \II{Pulslose elektrische Aktivität} (\I{PEA}): Das
    Herz reagiert nicht auf elektrische Reize
  \end{itemize}
\end{itemize}
Beim (S)AED übernimmt das Gerät die Beurteilung ob ein
elektrischer Herzrhythmus schockbar oder nicht schockbar
ist. 
%
Die Herzrhythmen werden im Kapitel
\FullPrettyRef{topic:herzrhythmusstoerungen}, 
genauer besprochen.
}{
\begin{itemize}
  \item schockbare Rhythmen
  \begin{itemize}
    \item Kammerflimmern
    \item pulslose ventrikuläre Tachykardie
    \end{itemize}
  \item nicht-schockbare Rhythmen
  \begin{itemize}
    \item Asystolie
    \item pulslose elektrische Aktivität
    \end{itemize}
\end{itemize}
}{}{}{}



\begin{figure}[h]
\caption{Sinusrhythmus und reanimationspflichtige Rhythmen}

\subfloat[\Z{(Langsamer) Sinusrhythmus}
\SourcePicture{WM/PD}] {\includegraphics[width=\linewidth]{./images/wikipedia-pd/Lead_II_rhythm_generated_sinus_bradycardia}}

\subfloat[\Z{Reanimationspflichtig, schockbar: Ventrikuläre Tachykardie} \SourcePicture{WM/PD}]
{\includegraphics[width=\linewidth]{./images/wikipedia-pd/Lead_II_rhythm_ventricular_tachycardia_Vtach_VT}}

\subfloat[\Z{Reanimationspflichtig, schockbar: Kammerflimmern}
\SourcePicture{WM/PD}]
{\includegraphics[width=\linewidth]{./images/wikipedia-pd/Lead_II_rhythm_generated_ventricular_fibrilation_VF}}

\subfloat[\Z{Reanimationspflichtig, nicht schockbar: Asystolie}
\SourcePicture{WM/PD}]
{\includegraphics[width=\linewidth]{./images/wikipedia-pd/Lead_II_rhythm_generated_asystole}}
\end{figure}

\clearpage % TODO temporary clearpage
\TextSlide{Verwendung eines manuellen Defibrillators}{
Grundsätzlich ist jedes Gerät so zu verwenden wie es der
Hersteller in der \EE{Bedienungsanleitung} vorschreibt. 


Manuelle Defibrillatoren sind meist Bestandteil eines Multifunktionsgerätes, welches oft auch umfangreiche Überwachungsmöglichkeiten bietet. 

\begin{description}
  \item[Therapieelektroden] 
  Bei aktuellen Geräten werden für die Defibrillation Klebeelektroden für die Defibrillation eingesetzt, wie sie auch bei (S)AEDs zum Einsatz kommen. 
  %
  Alternativ, oder bei älteren Geräten auch ausschließlich, können auch \II{Hardpaddles} verwendet werden. 
  %
  Diese bestehen aus 2 Elektroden als Aufsatzfläche und jeweils einem isolierten Haltegriff. Die Positionierung der Elektroden am Körper des Patienten erfolgt wie bei den Klebeelektroden.
  %
  Zusätzlich wird noch ein \I{Elektrodengel} benötigt, 
  um den Leitungswiderstand herabzusetzen. 
  Anders als in Filmen dargestellt, 
  dürfen die Elektroden der Paddles keinesfalls gegeneinander gerieben werden um das Elektrodengel zu verreiben!
  %
  Ausserdem muss bei Verwendung von Paddles ein \E{Anpressdruck} von mind. 8\Kg* angewendet werden.
  %
  
  Die Verwendung von Hardpaddles führt zu schlechteren Ergebnissen 
  und soll daher nur wenn nicht anders möglich erfolgen 
  \cite{Erc-2015-Section-03,Erc-2015-Sektion-03-De}. 
  \item[Manueller Modus] 
  Viele manuelle Defibrillatoren besitzen einen (S)AED-Modus, 
  welcher oft bei Starten des Geräts automatisch aktiviert ist.
  %
  Bei Bedienung durch Personal, 
  welches das abgeleitete EKG interpretieren kann und darf, 
  ist es sinnvoll, 
  in den manuellen Modus zu schalten, 
  um die Zeit für die Rhythmusanalyse möglichst kurz zu halten.
  \item[Analyse] 
  Die Analyse erfolgt im manuellen Modus durch einen kompetenten 
  und dazu befugten Helfer 
  zum gleichen Zeitpunkt wie im (S)AED-Modus,
  also alle 2\Min* 
  bzw. 5 Zyklen.
  \item[Auswahl der Energie] 
  Im manuellen Modus ist die Auswahl der abzugebenden Energie möglich. 
  Bei Erwachsenen ist bei der Defibrillation der Maximalwert zu verwenden, dieser ist je nach Gerät und Impulsform unterschiedlich. Bei Kindern beträgt die Energie 4 Joule pro {\Kg}.
  \item[Synchronisation] 
  Bei vorhandenen Kammerkomplexen im EKG kann eine Schockabgabe zu bestimmten Zeiten des Herzzyklus (\enquote*{vulnerable Phase}) Kammerflimmern ausgelöst werden. 
  %
  Um dies zu verhindern kann die Schockabgabe synchronisiert mit dem Herzzyklus erfolgen, der Defibrillator gibt dann denn Stromstoß in einem gefahrlosen Intervall ab.
  %
  Neuere Geräte verfügen oft über eine Rhythmuserkennung und \enquote*{\I{Auto-Sync}}-Funktion, hier entfällt die manuelle Einstellung.
  %
  \begin{Danger}
    \item Wird bei einem irregulären Rhythmus 
    (Kammerflimmern) 
    die Sync-Funktion aktiviert, 
    kann das Gerät keine Kammerkomplexe erkennen 
    und es erfolgt \E{keine} Schockabgabe!
  \end{Danger}
  \item[Schockabgabe] Vor der Schockabgabe hat sich der Helfer zu versichern, dass die Umgebung und das Personal sicher ist und niemand den Patienten berührt. Er hat dies auch laut kundzutun:
  \begin{enumerate}
    \item Kontrollblick auf das EKG: Hat sich der Rhythmus geändert, ist der Rhythmus defibrillierbar?
    \item \Speech{\enquote{Alle weg vom Patienten, wir schocken!}}
    \item Sicherungsblick: Berührt keiner den Patienten? Ist die Umgebung sicher (Nässe, Dämpfe, Sauerstoff vom Patienten entfernt, \ldots)?
    \item \Speech{\enquote{Achtung! Schock wird ausgelöst!}}
    \item Auslösen des Schocks.
  \end{enumerate}
\end{description}

%Bei
%der Verwendung von AEDs kommt es meist zu folgendem
%allgemeinen Vorgehen:
%
%Das Gerät wird in unmittelbarer Reichweite zu Patient und
%Helfer positioniert, die Gerätetasche wird geöffnet. Die
%Klebeelektroden sind Einmalprodukte und daher zusätzlich
%verpackt. Nach Öffnen der Schutzhülle werden die Kabel der
%Elektroden an das Gerät angeschlossen. Danach wird einzeln
%von jeder der beiden Elektroden die Schutzfolie abgezogen
%und die Elektrode \E{faltenfrei} auf den Patienten geklebt.
%Ist die Schutzfolie einmal abgezogen darf die Elektrode
%nicht mehr abgelegt werden, da die Klebeschicht sonst
%verschmutzt würde und wirkungslos wird. Ggf. müssen die
%Brusthaare des Patienten vor dem Aufkleben abrasiert werden.
%Ein Einmalrasierer ist i.\,d.\,R. in der Gerätetasche
%vorrätig.
%Schlecht geklebte Elektroden (Falten) oder nicht
%abrasiertes, dichtes Brusthaar führen zu einem unnötig hohen
%elektrischen Widerstand im Stromfluss und machen die
%Elektroschocks daher ineffizient. Die \EE{Positionierung der
%Elektroden} ist in Abbildung \ref{fig:defibrillation}
%ersichtlich. Durch diese Positionierung ist sicher gestellt,
%dass ein Großteil des elektrischen Feldes zwischen den
%Elektroden durch das Herz läuft.
%
%Sind die Elektroden angeklebt, kann mit der Beurteilung des
%elektrischen Herzrhythmus begonnen werden. Die
%\EE{Rhythmusanalyse} (kurz Analyse) führt das Gerät auf
%Knopfdruck selbständig durch. Es ertönt eine Warnung, z.\,B.
%\Speech{\enquote{Analyse läuft! Patienten nicht berühren!}}
%Der Patient (oder elektrisch leitende Gegenstände (z.\,B.
%Metalle, Wasser) die mit dem Patient in direktem Kontakt
%stehen) dürfen vom Helfer nicht berührt werden. Nach der
%Analyse wird den Helfern per Sprachausgabe mitgeteilt, ob
%ein schockbarer oder ein nicht-schockbarer Rhythmus
%vorliegt. Wird ein Schock empfohlen, kann mittels
%Tastendruck der Elektroschock ausgelöst werden.
%Selbstverständlich darf auch während der \EE{Schockabgabe}
%kein Kontakt zum Patienten besteht.
%\EE{Helfer müssen hier sorgfältig auf Sicherheit achten
%(warnen $\rightarrow$ zurücktreten $\rightarrow$
%kontrollieren)!}

}{
\begin{itemize}
  \item Verwenden gem. Bedienungsanleitung!
  \item Positionierung
  \item Aufkleben der Elektroden
  \item Manueller Modus
  \item Anaylse: 
  Alle 2\Min* bzw. 5 Zyklen durch kompetenten und befugten Helfer
  \item Energieauswahl: 
  
  Erwachsene Maximum, 
  
  Kinder 4\,J\,/\,Kg
  \item Synchronisation (\enquote*{Sync}): 
  Bei Vorhandenen Kammerkomplexen
  \item Schockabgabe
  \item Sicherheitshinweise
\end{itemize}
}{}{}{}

\bigskip

\begin{BoxRed}[title={\flushleft Sicherheits- und weitere Bedienungshinweise}]
\begin{itemize}
  \item \StepHeading{Fehlströme}: Der Defibrillator darf nicht in einer \EE{nassen
  Umgebung} oder auf \EE{Metallunterlagen} eingesetzt werden,
  da Wasser und die meisten Metalle den elektrischen Strom
  leiten.
  Dadurch sind Helfer bei der Schockabgabe gefährdet. 
  Patienten, die aus dem Wasser gerettet werden,
  müssen daher vor der Schockabgabe abgetrocknet und auf
  trockenen Untergrund gebracht werden.
  \item \StepHeading{Brandgefahr}: Der Defibrillator darf nicht 
  in der Nähe von entzündlichen Substanzen 
  oder in Umgebungen, 
  in welchen \EE{Explosionsgefahr} herrscht,
  eingesetzt werden!
  \item \StepHeading{Kinder}: 
  Viele Geräte dürfen bei {Kindern}
   erst ab einem 
  bestimmten Körpergewicht oder Alter 
  (je nach Herstellerangabe)
  eingesetzt werden.
  \item \StepHeading{Sauerstoffquellen}: 
  Während der Schockabgabe sind etwaige
  \E{Sauerstoffquellen} so weit wie möglich von den
  Elektroden weg zu halten (Richtwert ca. 1\Meter*).
%   \item \E{Unterkühlte Patienten} mit einer
%   Körperkerntemperatur unter 30\DegC{} dürfen ebenfalls
%   nicht defibrilliert werden. % ex mit ERC 2010
  \item \StepHeading{Interferenzen}: Folgende Umstände können die Beurteilung des Herzrhythmus verfälschen 
  \begin{itemize}
    \item \StepHeading{Bewegungsartefakte}:Eigenbewegung des Patienten,
    Herzdruckmassage,
      Manipulationen am Patienten,
      Manipulationen am Defibrillator oder Kabel, maschinelle
      Beatmung, Transport, sonstige Erschütterungen. 
    \item \StepHeading{Elektromagnetische Felder}: 
    Bei Starkstromleitungen muss ein Abstand von
    mindestens 3\Meter* gehalten werden.
  \end{itemize}
  \item \StepHeading{Medikamentenpflaster} müssen vor der
  Defibrillation entfernt werden.
  \item Die Klebeelektroden müssen \E{ohne Lufteinschlüsse}
  aufgeklebt werden.
  \item \StepHeading{Paddles}: Bei Verwendung von Paddles muss Elektrodengel und ein Anpressdruck von mind. 8\Kg* verwendet werden
\end{itemize}
Zusätzliche Sicherheits- und Bedienhinweise sind in der
Bedienungsanleitung des jeweiligen Geräts nachzulesen!

\begin{center}
  \Ca{Lebensgefahr!} 
\end{center}

\end{BoxRed}

\Picture
{Ein Multifunktionsgerät mit integriertem Defibrillator (Corpuls\textsuperscript{3}\texttrademark)}% Titel
{}% Beschreibung
{}% Label
{Christoph Pallinger}% Quelle
{MfG}% Lizenz
{./images/pallinger-christoph-aass/Corpuls_32863-00943pt}% Datei
{}% Breitenmodifizierer
{}% Float


% \TakeNotesLarge

\subsection{\HeadingKeyword{Airwaymanagement}}

\Text
{Einleitung}%1 Überschrift
{%
Ziel des Atemwegsmanagements (Airwaymanagement) ist die ausreichende Belüftung (Ventilation) der (unteren) Atemwege.
%
Dazu muss einerseits der Atemweg freigehalten werden
und andererseits einer möglichen Aspiration von Mageninhalt, Blut etc. vorgebeugt werden. 
%
Zum Erreichen dieser Ziele gibt es verschiedene Techniken,
welche sich bezüglich Aufwand und Komplexität deutlich unterscheiden.
%
Es gibt nicht genügend Beweise
für oder gegen eine spezielle Technik zum
\I{Airwaymanagement} beim Kreislaufstillstand, 
vielmehr ist die Auswahl von 
der Situation, 
dem verfügbaren Material 
sowie dem Personal 
und dessen Ausbildung, 
Training 
und Routine abhängig.
}%2 Text


\TextSlidePicture
{Endotracheale Intubation}%1 Titel
{%
\DefinitionInPara{Als
\T[Intubation!endotracheale]{endotracheale Intubation}
(\T{ETI}) wird das Einführen eines Tubus in die Trachea
bezeichnet.} Sie ist eine Maßnahme des invasiven
Atemwegsmanagements und wird als Goldstandard hinsichtlich
der Qualität der Atemwegssicherung angesehen, dies
berücksichtigt aber nicht die Sicherheit der Anwendung.

Der Tubus wird mithilfe eines Laryngoskops unter Darstellung
der Stimmritze unter Sicht in die Luftröhre eingeführt.
An den gängigen Tuben ist ein Cuff angebracht, welcher ein
Abdichten des Atemweges ermöglicht\footnote{Manche Tuben,
z.\,B. für Klinkinder haben keinen Cuff.}.

\subparagraph{Kontrolle der Tubuslage}
Das korrekte Platzieren des Tubus ist einer der kritischten
Punkte bei der endotrachealen Intubation. Typische Fehler
sind die Intubation der Speiseröhre (\E{ösophageale}
Intubation) oder die Intubation eines Hauptbronchus
(\E{bronchiale} Intubation bei zu tief vorgeschobenen Tubus, 
es wird dann nur \E{ein} Lungenflügel belüftet).

Grundsätzlich gibt es präklinisch drei Möglichkeiten, die
korrekte Lage des Tubus zu überprüfen:
\begin{enumerate}
  \item Visuelle Darstellung mittels Laryngoskop
  \item Auskultation von Magen und Lungenflügel
  \item Messung der \ce{CO2}-Abatmung mittels Kapnograhie
\end{enumerate}
Es sollen immer \EE{alle drei
Möglichkeiten} angewendet werden.
Es wird jedoch zunehmend ein größerer Wert auf die
\EE{Kapnographie} gelegt, um die Tubuslage zu kontrollieren,
die CPR-Qualität abzuschätzen und eine frühzeitige Erkennung
der Rückkehr des spontanen Kreislaufs zu ermöglichen. Die
Kapnographie ist die verlässlichste Methode zur
Lagekontrolle \cite{Grmec2002}. Besonders
geeignet sind dazu Geräte, welche die \ce{CO2}-Abatmung als Kurve darstellen. Wenn nach 6
Beatmungen \ce{CO2} nachweisbar ist, kann man davon
ausgehen, dass der Tubus in den unteren Atemwege zu liegen
gekommen ist. Die Kapnographie kann jedoch \E{keine} Aussage
darüber treffen, ob die \E{Trachea oder nur ein einzelner
Bronchus} intubiert wurde.

Die Kapnographie ist Voraussetzung 
für die Durchführung einer endotrachealen Intubation. 
Wenn kein Kapnograph mit Kurvendarstellung vorhanden ist, 
sollen alternative Atemwegshilfsmittel zum Einsatz kommen
%
\ZFootnote{\Speech{\enquote{Anyone 
attempting tracheal intubation must be well trained
and equipped 
with waveform capnography. 
%
In the absence of these prerequisites, 
consider use of bag-mask ventilation 
and/or an SGA
until appropriately experience 
and equipped personnel are present.
}} \cite{Erc-2015-Section-04}.
\Commentary[Airwaymanagement / Kapnographie]{Zur
Problematik der Verfügbarkeit der Kapnographie siehe
auch 
\cite{Genzwuerker2007}
(\fullcite{Genzwuerker2007}) 
\cite{Timmermann2007}
und darin enthaltene
Referenzen.
}%Commentary
}%ZFootnote
.
}%2 Text
{%
\begin{itemize}
  \item Tubus in Trachea
  \item Laryngoskop
  \item Kontrolle Tubuslage:
  \begin{itemize}
    \item Visuelle Darstellung
    \item Auskultation Magen, Lunge
    \item Kapnographie (\ce{CO2})
  \end{itemize}
\end{itemize}
}%3 Slide
{./images/motal-aass/intubation/dsc_4837-AASS-0030mm}%4 Bilddatei
{Endotracheale Intubation}%5 Bildtitel
{}%6 Bildbeschreibung
{Michael Motal}%7 Quelle
{MfG}%8 Lizenz

\begin{Synopsis}
  \SynopsisItem{Die Kapnographie ist Voraussetzung für die endotracheale Intubation.}
\end{Synopsis}



\Picture
[]% Float
{\ce{CO2}-Kurve}% Titel
{Die Kapnometrie ist wichtig, um die korrekte Lage des Tubus in der Luftröhre zu bestätigen.}% Beschreibung
{}% Label
{Sebastian Gabriel}% Quelle
{}% Lizenz
{./images/gabriel-sebastian-aass/Monitor-Beatmungseinstellungen-Live-00943pt-4}% Datei
{}% Breitenmodifizierer

\TextSlidePicture
{Supraglottische Atemwegshilfsmittel}%1 Titel
{%
In den letzten Jahren gewinnen die
supraglottischen\footnote{\Lang{lat.} supra: oberhalb, über;
\Lang{Term.} glottis: Der aus den Stimmfalten bestehende
Teil des Kehlkopfinnenraums.} Hilfsmittel aufgrund ihrer
verhältnismäßig einfachen Anwendbarkeit immer mehr an
Bedeutung. Schon die ERC-Leitlinien 2010 sahen die
supraglottischen Atemwegshilfen bei der Reanimation als
weitgehend gleichwertig zur klassischen endotrachealen
Intubation an
\cite[p.1319,1321]{Erc2010Sec4}.\footnote{Supraglottische
Methoden zur Atemwegssicherung werden oft als
\enquote*{alternativer Atemweg} bezeichnet, da sie oft eine
Alternative zur endotrachealen Intubation (\I{ETI})
angesehen werden. Die ETI gilt weithin als Goldstandard der
Atemwegssicherung, dennoch hat sie einige gewichtige
Nachteile und ist in vielen Situationen \E{nicht die Methode
der Wahl}. }

Bei den supraglottischen Atemwegshilfen wird der Schlauch
nicht in die Luftröhre eingeführt, sondern endet im Rachen
beziehungsweise oberhalb des Kehldeckels (supraglottisch).
Es gibt verschiedene Produkte, welche dieses Ziel auf
unterschiedliche Weise erreichen. Dazu gehören der
Larynxtubus{\texttrademark}, die Larynxmaske und der Combitubus.

Der Larynxtubus{\texttrademark} ist unter \FullPrettyRef{topic:larynxtubus} beschrieben.
}%2 Text
{}%3 Slide
{./images/pallinger-christoph-aass/Laringstubus_32884-00441pt}%4 Bilddatei
{Ein Larynxtubus}%5 Bildtitel
{}%6 Bildbeschreibung
{Christoph Pallinger}%7 Quelle
{MfG}%8 Lizenz




\begin{PictureSeries}{}{Übersicht: Lage Larynxtubus und endotrachealer Tubus
\SourcePicture{Hirtler}}
\PictureSeriesRow{2}
{./images/hirtler-aass/LageLarynxTubus}
{Im Vergleich: Lage eines Larynx-Tubus \ldots} 
{./images/hirtler-aass/LageTubus-color}
{\ldots\ und eines endotrachealen Tubus.}
{./images/hirtler-aass/PositionGuedel-color}
{Position des Guedel-Tubus im Querschnitt}
{empty}
{}
\end{PictureSeries}


\Text
{Streitfrage: Endotracheale Intubation oder Alternative?}
{%
Die endotracheale Intubation (ETI) ist grundsätzlich eine
sehr anspruchsvolle Tätigkeit.  Es gibt viele
Fehlerquellen, beispielhaft seien die Intubation des
Ösophagus oder nur eines Bronchus, Verletzung
beim Intubationsvorgang oder ein Bronchspasmus genannt.
Zahlreiche Studien haben eine hohe Inzidenz von
Fehlintubationen beschrieben.
Speziell im 
außerklinischen Bereich erhöht sich der
Schwierigkeitsfaktor aufgrund der Gegebenheiten
(Platzangebot, vorhandenes Material, nicht nüchterne bzw.
nicht vorbereitete Patienten, Verletzungen,
Begleiterkrankungen, \ldots) deutlich
\cite{Timmermann2006,Thierbach2004}.

\Z{%
In diversen Studien wurden die Erfolgs- und
Komplikationsraten der ETI untersucht, dabei zeigen sich
wenig zufriedenstellende bis alarmierende Ergebnisse %
\cite{Nolan2001,
% Prehospital and resuscitative airway care: should the gold standard be reassessed?
% In the context of prehospital care and resuscitation,
% tracheal intubation has been regarded as the standard in
% airway treatment. The evidence for this status is rather
% weak. It does not take into account the level of training
% and experience of the personnel attempting intubation, and
% whether they use neuromuscular blockers. In unskilled
% hands, attempted tracheal intubation is harmful;
% unrecognized esophageal intubation is disastrous. When
% healthcare providers lack adequate skills in tracheal
% intubation, alternative airway devices, such as the
% laryngeal mask airway or the Combitube, may be better
% options than a simple facemask. Healthcare personnel using
% any of these devices should be adequately trained and
% maintain frequent practice.
Wang2011,% Out-of-hospital airway management in the
% United States.In this study characterizing
% out-of-hospital airway management across the United
% States, we observed low out-of-hospital ETI success rates.
Cobas2009, % Prehospital intubations and mortality: a level
% 1 trauma center perspective.
% This prospective study showed a 31\% incidence of failed
% PHI in a large metropolitan trauma center. We found no
% difference in mortality between patients who were properly
% intubated and those who were not, supporting the use of
% bag-valve-mask as an adequate method of airway management
% for critically ill trauma patients in whom intubation
% cannot be achieved promptly in the prehospital setting.
Jones2004, % Emergency physician-verified out-of-hospital
% intubation: miss rates by paramedics. 
% The rate of unrecognized, misplaced out-of-hospital
% intubations in this urban, midwestern setting was 5.8\%
Wang2006, % Paramedic intubation errors: isolated events or
% symptoms of larger problems?
% In our study population, we found that errors occurred in
% 22 percent of intubation attempts, with a frequency of up
% to 40 percent in selected ambulance systems. These
% findings indicate frequent errors associated with this
% life-saving technique. These events might be emblematic of
% larger issues in the structure and delivery of
% out-of-hospital emergency care.
Wang2006a, % Out-of-hospital endotracheal intubation: where
% are we?
% These studies highlight our limited understanding of
% out-of-hospital endotracheal intubation and the need for
% new strategies to improve airway support in the
% out-of-hospital setting.
Wang2008,
Stiell2004,
Stiell2008,
}. Die wichtigsten Ergebnisse dieser Studien sind: %
}% Z
\begin{itemize}
  \item Die Gefahr von Fehlintubationen, insbesonders von
  ösophagealen Fehllagen, ist hoch.
  Diese werden oft spät oder gar nicht
  erkannt (Risiko der unerkannten
  Tubusfehllage bei Patienten mit außerklinischem
  Kreislaufstillstand: \Range{0,5}{17}\PC*%
  \cite{Grmec2002, % 0,5%: Comparison of three different
  % methods to confirm tracheal tube placement in emergency intubation.
  % The esophageal intubations were followed by successful
  % endotracheal intubations without complications. The
  % capnometry (sensitivity and specificity 100\%) and
  % capnography (sensitivity and specificity 100\%) were
  % better than auscultation (sensitivity 94\% and specificity
  % 83\%) in confirming endotracheal tube placement in
  % non-arrest patients ( p<0.05). Capnometry was highly
  % specific (100\%) but not sensitive (88\%) for correct
  % endotracheal intubation in patients with cardiopulmonary
  % arrest (capnometry versus auscultation and capnometry
  % versus capnography, p<0.05).Capnography is the most
  % reliable method to confirm endotracheal tube placement in
  % emergency conditions in the prehospital setting.
  Lyon2010, % 2,4%: Field intubation of cardiac
  % arrest patients: a dying art?
  % The aetiology of cardiac arrest was medical in 95.2\%,
  % traumatic in 3.9\% and unrecorded in 0.9\%. Prehospital
  % intubation was attempted in 628 patients. Prehospital
  % intubation was successful in 573 patients. A significant
  % complication (multiple attempts, displaced endotracheal
  % tube or oesophageal intubation) occurred in 55 (8.8\%)
  % patients. 165 (20.8\%) patients survived to hospital
  % admission, of whom 110 had undergone prehospital
  % intubation. 55 patients who did not undergo prehospital
  % tracheal intubation survived to hospital admission.The
  % optimal method of maintaining an airway and ventilating an
  % OHCA patient has yet to be established. Prehospital
  % tracheal intubation for OHCA is associated with
  % significant complications and may reduce survival. The use
  % of tracheal intubation as a routine intervention should be
  % reconsidered. Ambulance services should consider adopting
  % alternative strategies in airway management.
  Jones2004, % 6%: Emergency physician-verified
  % out-of-hospital intubation: miss rates by paramedics.
  Pelucio1997, % Out-of-hospital experience with the syringe
  % esophageal detector device.
  % In the remaining 168 intubations, the EDD correctly
  % identified 5 of 10 esophageal intubations (sensitivity
  % 50\%; 95\% CI: 19\%, 81\%).
  % The EDD correctly identified 156 of 158 tracheal
  % intubations (specificity 99\%; 95\% CI: 96\%, 100\%).In
  % this paramedic field study, the EDD demonstrated poor
  % sensitivity for esophageal intubations.
  Jemmett2003, % 9%: Unrecognized misplacement of
  % endotracheal tubes in a mixed urban to rural emergency
  % medical services setting.
  % Of the studied patients, 12\% (13 of 109) were found to
  % have misplaced endotracheal tubes. For the patients with
  % unrecognized improperly placed tubes, 9\% (10 of 109) were
  % in the esophagus, 2\% (2 of 109) were in the right main
  % stem, and 1\% (1 of 109) were above the cords. Paramedics
  % serving urban and suburban areas did not perform
  % significantly better (p < 0.05) than intermediate-level
  % providers serving areas that are more rural.
  % The incidence of unrecognized misplacement of endotracheal
  % tubes by EMS personnel may be higher than most previous
  % studies, making regular EMS evaluation and the
  % out-of-hospital use of devices to confirm placement
  % imperative. The authors were unable to show a difference
  % in misplacement rates based on provider experience or
  % level of training.
  Katz2001, % Misplaced endotracheal tubes by paramedics
  % in an urban emergency medical services system.
  % On arrival in the ED, 25\% (27/108) of patients were found
  % to have improperly placed endotracheal tubes. Of the
  % misplaced tubes, 67\% (18/27) were found to be in the
  % esophagus, whereas in 33\% (9/27), the tip of the tube was
  % found to be in the hypopharynx, above the vocal cords. %
  Erc2010Sek04De,
  }
  \item Die Erfolgsraten bei wenig routinierten Anwendern
  sind schlecht\footnote{\Speech{\enquote{Die
  Misserfolgsraten der Intubation betragen in
  außerklinischen, wenig ausgelasteten Systemen
  mit Anwendern, die selten Intubationen durchführen, bis
  zu 50\PC*}}\cite{Erc2010Sek04De}.}.
  \item Die Probleme der ETI  betreffen grundsätzlich auch ärztliches
  Personal\footnote{Timmermann et.~al. 
  2007 \cite{Timmermann2007}: In dieser in Deutschland
  durchgeführten Studie wurde eine hohe Rate an
  Fehlintubationen (ösophageal, bronchial) durch Notärzte
  festgestellt.}. \cite{Timmermann2007,Grmec2002, % 0,5%:
  % Comparison of three different methods to confirm
  % tracheal tube placement in emergency intubation.
  }).
  \item Der Schulungsaufwand ist groß: 
  \Z{Zum Erlernen der
  Tätigkeit werden wenigstens 100 ETI gefordert, zur
  Aufrechterhaltung mindestens 10 ETI pro
  Jahr \cite{Timmermann2012}. Damit verbunden wäre eine
  Vollzeit-Trainingsphase von etwa einem halben Jahr
  \cite{Bernhard2012}.}
\end{itemize}

Es kommt hinzu, dass ein Vorteil für den Patienten
durch die ETI \E{nicht} nachgewiesen ist:
\Z{In 2008 fassten Nolan und Soar \cite{Nolan2008} zusammen, dass der Nutzen der ETI
bei der präklinischen Reanimation
nicht durch Studien gezeigt werden konnte, und dass einige
sogar das Gegenteil erkennen ließen
(\Speech{\enquote{Prehospital studies fail to show any
benefit from tracheal intubation during cardiopulmonary
resuscitation and many show harm.}}~\cite{Nolan2008}). Dies
deckt sich auch mit Beobachtungen aus anderen Studien \cite{
Stiell2004,%There was no improvement in the rate of survival
% with the use of advanced life support in any subgroup.
Cobas2009, % Prehospital intubations and mortality: a level
% 1 trauma center perspective. 
}.
Ein 2008 veröffentlichtes Cochrane-Review kam zu der
gleichen Schlussfolgerung, mit der Betonung, dass entsprechende, qualitative
hochwertige Studien fehlen \cite{Lecky2008}. 
} % Z


Die frühe endotracheale Intubation verliert somit an
Stellenwert.
Sie soll nur angewendet werden, wenn der Anwender hochqualifiziert
und sich sehr sicher ist,
bzw. regelmäßige, laufende Erfahrung
mit der Technik hat 
(\Speech{\enquote{It should be used
only when trained personnel are available to carry out the procedure with a
high level of skill and confidence. }} \cite{Erc2010Sec4,Erc-2015-Section-03}).
Ohne ausreichendes Training und Erfahrung wird die
Komplikationsrate als unakzeptabel hoch angesehen. Das
Aufschieben der endotrachealen Intubation bis zur Wiederkehr
des spontanen Kreislaufs, wird bereits in den ERC 2010-Leitlinien
als Alternative ausdrücklich genannt.
\cite{Nolan2001,
Nolan2008,
Lecky2008,
Erc2010Sec4}
Vor diesem Hintergrund bleibt festzuhalten, dass eine
unterlassene Intubation nicht -- wie oft behauptet -- mit
einer unterlassenen Hilfeleistung, sondern mit einer
unterlassenen Tradition gleichzusetzen ist.


Der traditionellen ETI steht eine stetige
Entwicklung an alternativen Airwaysystemen gegenüber, welche eine effektive und
technisch einfachere Sicherung des Atemweges ermöglichen
(supraglottsiche Atemwegshilfen).
Zwar müssen diese Systeme Abstriche hinsichtlich des
Aspirationsschutzes machen, wiegen dies jedoch durch ihre
einfach und sichere Anwendung auf %
\cite{Schalk2010, %
% Out-of-hospital airway management by paramedics and emergency physicians using laryngeal tubes.
% The LT-D/LTS-D represents a reliable tool for prehospital
% airway management in the hands of both paramedics and
% emergency physicians. It can be used as an initial tool to
% secure the airway until ETI is prepared, as a definitive
% airway by rescuers less experienced with ETI or as a
% rescue device when ETI has failed.
% The placement time was < or =45s (n=120), 46-90s (n=20)
% and >90s (n=7).
Wiese2009, % The use of the laryngeal tube disposable (LT-D)
% by paramedics during out-of-hospital resuscitation-an observational study concerning ERC guidelines 2005.
% The LT-D was used in 46\% of all cardiac arrest
% situations. Overall, the LT-D was successfully inserted in
% more than 90\% of all cases on first attempt. In 95\% of
% all cases, no problems concerning ventilation of the
% patient were described.As an alternative airway device
% recommended by the ERC in 2005, the LT-D may enable airway
% control rapidly and effectively. Additionally, by using
% the LT-D, a reduced "no-flow-time" and a better outcome
% may be possible.
}. 

}
{}
{}
{}
{}

\begin{Synopsis}
  \SynopsisItem{Wenn die erforderliche Routine und Sicherheit für die
  Durchführung einer ETI nicht gegeben ist 
  und die Indikation   für ein invasives Atemwegsmanagement besteht, 
  sollen \EE{in erster Linie supraglottische Atemwegshilfen}
  eingesetzt werden \cite{Timmermann2012}.}
\end{Synopsis}

\begin{Literary}
\fullcite{Timmermann2012}

\fullcite{Lecky2008}

\fullcite{Timmermann2007}

\fullcite{Erc-2015-Sektion-04-De}
\end{Literary}


%\subsubsection{\HeadingKeyword{Larynxtubus{\texttrademark}}}
%
%Siehe \FullPrettyRef{topic:larynxtubus}.



\subsection{\HeadingKeyword{Medikamente}}
\index{Adrenalin!Reanimation}%
\index{Amiodaron!Reanimation}%
\index{Atropin!Reanimation}%

\Text
{}%1 Überschrift
{%
Die Gabe von Medikamenten ist eine Routinemaßnahme während
der Reanimation. Typischerweise erfolgt die Medikamentengabe
über eine periphere Venenverweilkanüle
(Venflon{\texttrademark}, Braunüle{\texttrademark}).
Sinvoller weise wird standardmäßig auch eine Infusion mit
kristalloider Infusionslösung (Kochsalz-,
Ringerlösung, Ringerlaktat) angeschlossen, um eine schnelle
Spülung zu ermöglichen.

Sollte
die rasche Anlage einer solchen nicht gelingen, kann als
Alternative ein intraossärer Zugang mittels spezieller
Hilfsmittel gelegt werden.
Früher war es auch üblich, Medikamente über einen
Endotrachealtubus (\I{ET}) zu geben, dies wird nicht mehr
empfohlen.
}%2 Text



\TableWide
{Übersicht: Die wichtigsten bei der
Reanimation eingesetzten
Medikamente}%1 Titel
{}%2 Beschreibung
{}%3 Label/Hook
{}%4 Quelle
{}%5 Lizenz
{%
\begin{tabulary}{\linewidth}{lLLL}
\toprule\addlinespace
\noindent\TabHeading{Wirkstoff}
&
\noindent\TabHeading{Erläuterungen}
&
\noindent\TabHeading{Indikation}
&
\noindent\TabHeading{Spezialitäten}
\\\addlinespace\midrule\addlinespace
\noindent\TabSub{Adrenalin}
&
Vasokonstriktor
&
Schockbarer Rhythmus

Nicht-Schockbarer Rhythmus
&
\noindent\RaggedRight\Z{L-Adrenalin Fresenius{\texttrademark} 2\Mg*\,/\,20\Ml*

L-Adrenalin Fresenius{\texttrademark} 0,5\Mg*\,/\,5\Ml*

Suprarenin{\texttrademark} 1\Mg*\,/\,1\Ml*}
\\\addlinespace
\noindent\TabSub{Amiodaron}
&
Anti\-ar\-rhythmi\-kum

Kl. III
&
Schockbarer Rhythmus
&
\noindent\RaggedRight\Z{Sedacoron{\texttrademark} 150\Mg*\,/\,3\Ml*}
\\\addlinespace
\noindent\TabSub{\Z{Atropin}}
&
\noindent\Z{Anti\-para\-sym\-patho\-mi\-meti\-kum}
&
\noindent\Z{Nicht-Schockbarer Rhythmus (seit 2010 nicht mehr
Standard)}
&
\noindent\RaggedRight\Z{Atropinum sulfuricum Nycomed{\texttrademark} 0,5\Mg*\,/\,1\Ml*}
\\\addlinespace\bottomrule
\end{tabulary}
}%6 Tabelle
{}%7 Float

% \begin{PictureSeriesWide}{}{Medikamente bei der Reanimation (Amiodaron nicht abgebildet) \SourcePicture{ASBÖ/gjh}}
% \PictureSeriesRowWide{3}
% {./images/ASBOE-Grassl-gjh-1/Medikamente0395_0197-00441pt}
% {L-Adrenalin \enquote*{Fresenius\texttrademark} 2\Mg* 20\Ml* Ampulle.}
% {./images/ASBOE-Grassl-gjh-1/suprarenin-freigestellt-borderadded-00441pt}
% {Suprarenin{\texttrademark}-Ampulle
% (1\Mg* Adrenalin in 1\Ml*, dieses Medikament wird vor Gebrauch verdünnt)}
% {./images/ASBOE-Grassl-gjh-1/Medikamente0465_0267-00441pt}
% {Eventuell ultima ratio zur Behandlung reversibler Ursachen: Lyse}
% {}
% {}
% \end{PictureSeriesWide}


% TODO Grassl Foto: Reanimationsmedikamente



\TableWide
{Standard-Reanimationsmedikamente ERC {2010} (vereinfacht)}%1 Titel
{}%2 Beschreibung
{}%3 Label/Hook
{}%4 Quelle
{}%5 Lizenz
{%
\begin{tabularx}{\linewidth}{XXX}
\toprule\addlinespace
\noindent\TabHeading{Medikament}
&
\noindent\TabHeading{Schockbar}
&
\noindent\TabHeading{Nicht-schockbar}
\\\addlinespace\midrule\addlinespace
\noindent\TabSub{Adrenalin}
&
\EE{1\Mg*} \Abk{IV}

\E{nach} dem 3., 5., 7.,~{\ldots} Schock
&
\EE{1\Mg*} \Abk{IV}

ab Zugang, alle \VonBis{3}{5} Minuten
\\\addlinespace
\noindent\TabSub{Amiodaron}
&
\EE{300\Mg*} \Abk{IV}

einmalig \E{nach} dem {3.} Schock

Eventuell \EE{150\Mg*} \Abk{IV}

einmalig \E{nach} dem {5.} Schock
&
---
\\\addlinespace\bottomrule
\end{tabularx}
}%6 Tabelle
{}%7 Float


% \begin{PictureSeriesWide}{}{Es war einmal: Bekannte Medikamente, welcher nicht mehr Standard bei einer Reanimation sind \SourcePicture{ASBÖ/gjh}}
% \PictureSeriesRowWide{3}
% {./images/ASBOE-Grassl-gjh-1/Medikamente0426_0228-00441pt}
% {Nicht mehr Standard bei der Reanimation: Atropin}
% {./images/ASBOE-Grassl-gjh-1/Medikamente0398_0200-00441pt}
% {Auch nicht mehr Standard: Natrium-Bikarbonat (NaBic)}
% {empty}
% {}
% {}
% {}
% \end{PictureSeriesWide}


% TODO Grassl Foto: Ehemalige Reanimationsmedikamente

\subsection{\HeadingKeyword{Reversible} Ursachen}

\TextSlide
{4 H + HITS}%1 Überschrift
{%
Es gibt acht klassische Ursachen eines
Herz-Kreislaufstillstandes, welche grundsätzlich behebbar
(reversibel) sind:

\begin{multicols}{2}
\TableBaseModification
\begin{itemize}
  \item \EE{\Ca{H}ypoxie}

  (z.\,B. Atemwegsverlegung),
  \item \EE{\Ca{H}ypovolämie}

  (z.\,B. Blutung durch Trauma,   GI-Blutung, Aneurysmaruptur),
  \item \EE{\Ca{H}ypo-/Hyperkaliämie}

  (z.\,B. metabolische Störung),
  \item \EE{\Ca{H}ypothermie}

  (z.\,B. Ertrinkungsunfall)
  \item \EE{\Ca{H}erzbeuteltamponade}

  (z.\,B. Thoraxtrauma durch Stich oder Schuss)
  \item \EE{\Ca{I}ntoxikation}

  (z.\,B. Opiate, trizyklische Antidepressiva)
  \item \EE{\Ca{T}hromboembolie}

  (Lungenembolie, Herzinfarkt)
  \item \EE{\Ca{S}pannungspneumothorax}
\end{itemize}
\end{multicols}

Gemäß der Anfangsbuchstaben werden diese reversiblen Ursachen als 
\T{4 H + HITS} 
bezeichnet.
%
Leider lassen sich nicht alle dieser reversiblen Ursachen
präklinisch erkennen oder beheben, dennoch muss dies -- wenn
möglich -- versucht werden.
}%2 Text
{ %
\begin{itemize}
  \item \EE{\Ca{H}ypoxie}
  \item \EE{\Ca{H}ypovolämie}
  \item \EE{\Ca{H}ypo-/Hyperkaliämie}
  \item \EE{\Ca{H}ypothermie}
  \item \EE{\Ca{H}erzbeuteltamponade}
  \item \EE{\Ca{I}ntoxikation}
  \item \EE{\Ca{T}hromboembolie}
  \item \EE{\Ca{S}pannungspneumothorax}
\end{itemize}
}



\begin{PictureSeriesWide}{}{Bilderserie: Reanimation live}
\PictureSeriesRowWide
{2}% Anzahl Spalten
{./images/gabriel-sebastian-aass/EKG-Reanimation-VtBewegungsartefakt-00441pt}% Pfad Bild 1
{Monitor-Bild während einer Reanimation im AED-Modus: Ventrikuläre Tachykardie und Bewegungsartefakte durch die Herzdruckmassage.}% Text Bild 1
{./images/gabriel-sebastian-aass/Medikamente-LAdrenalin-Sammlung-Live-00441pt}% Pfad Bild 2
{Während einer Reanimation kann viel Adrenalin verbraucht werden.}% Text Bild 2
{./images/gabriel-sebastian-aass/}% Pfad Bild 3
{}% Text Bild 3
{}% Pfad Bild 4
{}% Text Bild 4
\end{PictureSeriesWide}


\section{Versorgung \HeadingKeyword{nach Reanimation}}

\Text
{Postreanimationsbehandlung}%1 Überschrift
{%
Wird bei der Reanimation der Kreislauf wiederhergestellt
(\T{ROSC}, \I{Resumption Of Spontaneous Circulation}) ist
die weitere Versorgung wichtig. 

Primär ist ein \EE{Einschätzungsblock} (erneut) durchzuführen. 
%
Bezüglich der Sofortmaßnahmen gelten sinngemäß die \EE{\TxMassVit}. 
%
Die Sauerstoffgabe hat einen Sp\ce{O2}-Zielbereich von \VonBis{94}{98}\PC*.
%
Ein zielgerichtetes \EE{Temperaturmanagement} mit einer Zielkörpertemperatur im Bereich \EE{\Range{32}{36}\DegC*} ist wichtig.
%
Eine höherer Körpertemperatur nach eienr Reanimation ist schädlich.

Die Ermittlung der \EE{Ursache}, 
sofern möglich, 
bzw. einer Verdachtsdiagnose 
ist wichtig um die weitere Versorgung zu planen.
%
Hevorzuheben ist die Notwendigkeit einer \EE{Herzkatheteruntersuchung} 
nach einem Kreislaufstillstand 
mit vermutet kardialer Ursache.

Dem Post-Cardiac-Arrest-Syndrome sowie dessen Behandlung
wird vermehrt Beachtung geschenkt.
Strukturierte post-Resuscitation-Protokolle erhöhen die
Überlebenswahrscheinlichtkeit und werden zunehmend
implementiert. 

\cite{Erc-2015-Sektion-05-De,Erc-2015-Section-05}

}%2 Text




\section{\HeadingKeyword{Abbruch} der Reanimation}

\Text
{}%1 Überschrift
{%
In ca. \VonBis{70}{95}\PC* der Fälle ist der
Reanimationsversuch erfolglos und der Patient verstirbt
\cite{StriebelAnaesthesie2-Band2,Erc2005DeSec8}.
Grundsätzlich wird eine Reanimation fortgesetzt, solange ein
Kammerflimmern oder ein ähnlicher Rhythmus vorliegt. Es ist
weitgehend akzeptiert, eine Reanimation abzubrechen, wenn
trotz aller ALS-Maßnahmen und Fehlen von reversiblen
Ursachen für 20 Minuten eine Asystolie besteht
\cite{Erc2010Sec10,Bonnin1993,StriebelAnaesthesie2-Band2}.
Gewisse Bedingungen, etwa eine Hypothermie zum Zeitpunkt
des Kreislaufstillstands, steigern die Chancen der Wieder-
herstellung ohne neurologische Schäden,
sodass normale Prognosekriterien (z.\,B.
eine über 20\Min* andauernde Asystolie)
nicht anwendbar sind. \autocite{Erc2005DeSec8}

Abgesehen von den rein medizinischen Faktoren,
welche vor allem den zu erwartenden Reanimationserfolg im
Blick haben, gibt es auch noch andere Umstände, welche
Einfluss auf die Entscheidung haben. Besonders
berücksichtigenswert ist der mutmaßliche oder vorab
geäußerte Patientenwille, z.\,B. im Rahmen einer
verbindlichen oder beachtlichen Patientenverfügung.

Die Abbruchsentscheidung ist grundsätzlich \EE{von einem
Arzt} zu treffen (sofern keine dringenden Gründe wie
unsichere Umgebung oder Erschöpfung vorliegen).
}%2 Text


\begin{Literary}
\begin{description}
  \item[\cite{www:Erc}] \fullcite{www:Erc}

  \item[\cite{www:ResuscitationJournal}] \fullcite{www:ResuscitationJournal}
  
\minisec{ERC-Guidelines, englische Originalfassung}

  \item[\cite{Erc-2015-Section-01}] \fullcite{Erc-2015-Section-01,}

  \item[\cite{Erc-2015-Section-02}] \fullcite{Erc-2015-Section-02}

  \item[\cite{Erc-2015-Section-03}] \fullcite{Erc-2015-Section-03}

  \item[\cite{Erc-2015-Section-04}] \fullcite{Erc-2015-Section-04}

  \item[\cite{Erc-2015-Section-05}] \fullcite{Erc-2015-Section-05}

  \item[\cite{Erc-2015-Section-06}] \fullcite{Erc-2015-Section-06}

  \item[\cite{Erc-2015-Section-07}] \fullcite{Erc-2015-Section-07}

  \item[\cite{Erc-2015-Section-08}] \fullcite{Erc-2015-Section-08}

  \item[\cite{Erc-2015-Section-09}] \fullcite{Erc-2015-Section-09}

  \item[\cite{Erc-2015-Section-10}] \fullcite{Erc-2015-Section-10}

  \item[\cite{Erc-2015-Section-11}] \fullcite{Erc-2015-Section-11}

\minisec{ERC-Guidelines, deutsche Übersetzung}

  \item[\cite{Erc-2015-Sektion-01-De}] \fullcite{Erc-2015-Sektion-01-De}

  \item[\cite{Erc-2015-Sektion-02-De}] \fullcite{Erc-2015-Sektion-02-De}

  \item[\cite{Erc-2015-Sektion-03-De}] \fullcite{Erc-2015-Sektion-03-De}

  \item[\cite{Erc-2015-Sektion-04-De}] \fullcite{Erc-2015-Sektion-04-De}

  \item[\cite{Erc-2015-Sektion-05-De}] \fullcite{Erc-2015-Sektion-05-De}

  \item[\cite{Erc-2015-Sektion-06-De}] \fullcite{Erc-2015-Sektion-06-De}

  \item[\cite{Erc-2015-Sektion-07-De}] \fullcite{Erc-2015-Sektion-07-De}

  \item[\cite{Erc-2015-Sektion-08-De}] \fullcite{Erc-2015-Sektion-08-De}

  \item[\cite{Erc-2015-Sektion-09-De}] \fullcite{Erc-2015-Sektion-09-De}

  \item[\cite{Erc-2015-Sektion-10-De}] \fullcite{Erc-2015-Sektion-10-De}

  \item[\cite{Erc-2015-Sektion-11-De}] \fullcite{Erc-2015-Sektion-11-De}
\end{description}
\end{Literary}

\PictureWide
	{ALS}%1
	{Der ALS-Algorithmus}%2
	{\label{fig:als}}%3
	{Emhofer, nach \cite{Erc-2015-Section-03}}%4
	{}%5
	{./images/emhofer-josef-aass/ReanimationAls-2015-3}%6
	{1}%7
%	{H}%8

\nocite{Uray2008}


\nocite{Erc2010Sec1}
\nocite{Erc2010Sec2}
\nocite{Erc2010Sec3}
\nocite{Erc2010Sec4}
\nocite{Erc2010Sec5}
\nocite{Erc2010Sec6}
\nocite{Erc2010Sec7}
\nocite{Erc2010Sec8}
\nocite{Erc2010Sec9}
\nocite{Erc2010Sec10}
\ChapterEnd
